
\documentclass[a4paper,10pt]{report}

%\usepackage[landscape]{geometry}
\usepackage[cm-default]{fontspec}
\setromanfont{FreeSerif}
\setsansfont{FreeSans}
\setmonofont{FreeMono}
\usepackage[utf8x]{inputenc}
\usepackage{fontspec}
\usepackage{xunicode}
\usepackage{xltxtra}
\usepackage{xgreek}
\usepackage{amsmath}
\usepackage{unicode-math}
\usepackage{ulem}
\usepackage{color}
\usepackage{verbatim}
\usepackage{nopageno}
\usepackage{graphicx}
\usepackage{textpos}
\setlength{\TPHorizModule}{1cm}
\setlength{\TPVertModule}{1cm}
%\usepackage[colorgrid,texcoord]{eso-pic}
\usepackage[outline]{contour}
\usepackage{wrapfig}
\usepackage{url}
\usepackage{color}
\usepackage{tikz}
\usepackage{fancybox,fancyhdr}
\usepackage{subfigure}
\usepackage{pstricks}
\usepackage{epsfig}
\usepackage{multicol}
\usepackage{listings}
\usepackage{enumerate}
\usepackage{hyperref}
\hypersetup{
  bookmarks=true,
  bookmarksopen=true,
  pdfborder=false,
  pdfpagemode=UseNone,
  raiselinks=true,
  pdfhighlight={/P},
  colorlinks,
  citecolor=black,
  filecolor=black,
  linkcolor=black,
  urlcolor=black
}


\usepackage{fontspec}
\usepackage{xunicode}
\usepackage{xltxtra}


% Margins
%\setlength{\textwidth}{24cm}

%\setlength{\voffset}{0in}
%\setlength{\textheight}{6.8in}
% Colors
\definecolor{cbrown}{rgb}{0.49,0.24,0.07}
\definecolor{cmpez}{rgb}{0.92,0.88,0.79}
\definecolor{cmpez2}{rgb}{0.62,0.33,0.34}
\definecolor{cwhite}{rgb}{1,1,1}
\definecolor{cblack}{rgb}{0,0,0}
\definecolor{cred}{rgb}{0.9,0.15,0.15}
\definecolor{clblue}{rgb}{0.57,0.85,0.97}
\definecolor{corange}{rgb}{0.86,0.49,0.18}
\definecolor{clorange}{rgb}{0.95,0.84,0.61}
\definecolor{cyellow}{rgb}{1,0.95,0.16}
\definecolor{cgreen}{rgb}{0.69,0.8,0.31}
\definecolor{clbrown}{rgb}{0.78,0.67,0.43}
\definecolor{cmagenta}{rgb}{0.79,0.05,0.54}
\definecolor{cgray}{rgb}{0.76,0.73,0.63}
\definecolor{bluesite}{rgb}{0.05,0.43,0.67}
\definecolor{redsite}{rgb}{0.90,0.15,0.15}
\definecolor{backsite}{rgb}{0.07,0.07,0.07}
% ----------- from Costas -----------------------------

%---------------   Main Fonts    -----------------%   

% \setmainfont[Mapping=tex-text]{Calibri}
% \setmainfont[Mapping=tex-text]{Times New Roman}
%\setmainfont[Mapping=tex-text]{Droid Serif}
%\setmainfont[Mapping=tex-text]{Cambria}
%\setmainfont[Mapping=tex-text]{cm-unicode}
% \setmainfont[Mapping=tex-text]{Gentium}
% \setmainfont[Mapping=tex-text]{GFS Didot}
% \setmainfont[Mapping=tex-text]{Comic Sans MS}
 \setmainfont[Mapping=tex-text]{Ubuntu}
% \setmainfont[Mapping=tex-text]{Myriad Pro}
% \setmainfont[Mapping=tex-text]{CMU Concrete}
%\setmainfont[Mapping=tex-text]{DejaVu Sans}
%\setmainfont[Mapping=tex-text]{KerkisSans}
%\setmainfont[Mapping=tex-text]{KerkisCaligraphic}
% \setmainfont[Mapping=tex-text]{Segoe Print}
% \setmainfont[Mapping=tex-text]{Gabriola}

%------------------------------------------------%

% ------------ Mathematics Fonts ----------------%
\setmathfont{Asana-Math.ttf}
% -----------------------------------------------%
% Misc
\author{Κασωτάκη Ε. - Σμαραγδάκης Κ.}
\title{www.math24.gr}
% ---------------- Formation --------------------%


\setlength\topmargin{-0.7cm}
\addtolength\topmargin{-\headheight}
\addtolength\topmargin{-\headsep}
\setlength\textheight{26cm}
\setlength\oddsidemargin{-0.54cm}
\setlength\evensidemargin{-0.54cm}
%\setlength\marginparwidth{1.5in}
\setlength\textwidth{17cm}

\RequirePackage[avantgarde]{quotchap}
\renewcommand\chapterheadstartvskip{\vspace*{0\baselineskip}}
\RequirePackage[calcwidth]{titlesec}
\titleformat{\section}[hang]{\bfseries}
{\Large\thesection}{12pt}{\Large}[{\titlerule[0.9pt]}]
%--------------------------------------------------%
\usepackage{draftwatermark}
\SetWatermarkText{www.math24.gr}
\SetWatermarkLightness{0.9}
\SetWatermarkScale{0.6}
%--------------------------------------------------%


\usepackage{xcolor}
\usepackage{amsthm}
\usepackage{framed}
\usepackage{parskip}

\colorlet{shadecolor}{bluesite!20}

\newtheorem*{orismos}{Ορισμός}

\newenvironment{notation}
  {\begin{shaded}\begin{theorem}}
  {\end{theorem}\end{shaded}}




% Document begins
\begin{document}
\pagestyle{fancy}
\fancyhead{}
\fancyfoot{}
\renewcommand{\headrulewidth}{0pt}
\renewcommand{\footrulewidth}{0pt}


\fancyhead[LO,LE]{
 %\textblockcolor{backsite}
 %\begin{textblock}{5}(-2,-0.55)
  %\rule{0cm}{1cm}
 %\end{textblock}
 \textblockcolor{bluesite}
 \begin{textblock}{5}(-1.5,-0.55)
  \rule{0cm}{1cm}
 \end{textblock}
 %\textblockcolor{bluesite}
 %\begin{textblock}{14}(5.5,-0.55)
 % \rule{0cm}{1cm}
 %\end{textblock}
 \begin{textblock}{0}(-1,-0.25)
 \color{cwhite} \begin{Large}www.math24.gr\end{Large}
 \end{textblock}
\begin{textblock}{0}(16,-0.4)
 \color{cwhite} \includegraphics[height=1.5cm]{math24_logo.png}
 \end{textblock}
\textblockcolor{white}
\begin{textblock}{17}(-1,28)
 \color{backsite} \begin{small}Copyright \textcopyright 2012, Κασωτάκη Ε.(ikasotaki@gmail.com) - Σμαραγδάκης Κ.(kesmarag@gmail.com)\end{small}
 \end{textblock}
}
 
\begin{shaded}
\begin{center}
\huge \textbf{Φυλλάδιο Ασκήσεων}\\
%Πρόσθεση ρητών αριθμών
\end{center} 
\textbf{Μαθηματικά Γ' Γυμνασίου} \hfill \textbf{Ημερομηνία Παράδοσης : \hspace{2em} }
\subsection*{Ονοματεπώνυμο :\hfill  \hspace{5em}}
\end{shaded}
\vspace{2em}
\begin{itemize}
 \item Δυνάμεις πραγματικών αριθμών
 \item Ιδιότητες δυνάμεων
\end{itemize}


\section*{Θεωρία - Ορισμός Δύναμης \hfill \small{}}
\textbf{Δύναμη} με βάση έναν πραγματικό αριθμό $α$ και εκθέτη ένα φυσικό αριθμό $ν\ge 2$ 
είναι το γινόμενο $ν$ παραγόντων ίσων με $α$ και συμβολίζεται με $α^{ν}$.\\
\textbf{π.χ} δύναμη με βάση $3$ και εκθέτη $4$: $3^{4}=3\cdot3\cdot3\cdot3$

\section*{Θεωρία - Ιδιότητες Δυνάμεων \hfill \small{}}
\begin{itemize}
 \item $α^{μ}\cdot α^{ν}=α^{μ+ν}$\\ 
       π.χ $3^{2}\cdot3^{4}=3^{2+4}=3^{6}$ \\
 \item $α^{μ}: α^{ν}=α^{μ-ν}$\\ 
       π.χ $2^{5}:2^{3}=2^{5-3}=2^{2}$\\
       π.χ $6^{5}:6^{11}=6^{5-11}=6^{-6}$ \\
 \item $(αβ)^{ν}=α^{ν}β^{ν}$\\
       π.χ $(3\cdot6)^{2}=3^{2}\cdot^{2}$\\
 \item $\bigg(\dfrac{α}{β}\biggl)^{ν}=\dfrac{α^{ν}}{β^{ν}}$\\
       π.χ $\bigg(\dfrac{2}{7}\biggl)^{3}=\dfrac{2^{3}}{2^{7}}$\\
 \item $(α^{μ})^{ν}=α^{μν}$\\
       π.χ $(7^{2})^{-3}=7^{2\cdot(-3)}=7^{-6}$\\
 \item $\bigg(\dfrac{α}{β}\biggl)^{-ν}=\bigg(\dfrac{β}{α}\biggl)^{ν}$\\
       $\bigg(\dfrac{3}{11}\biggl)^{-4}=\bigg(\dfrac{11}{3}\biggl)^{4}$\\
\end{itemize}

\newpage

\section*{Άσκηση 1  \hfill \small{10 μονάδες}}
Να αντιστοιχίσετε κάθε στοιχείο της αριστερής στήλης με ένα στοιχείο της δεξιάς στήλης
\begin{itemize}
\begin{multicols}{2}
 \item $2^{3}\cdot2^{7}$
 \item $2^{3}:2^{7}$
 \item $(2\cdot3)^{3}$
 \item $\bigg( \dfrac{2}{3}  \biggr)^{3}$
 \item $(2^{3})^{2}$
 \item $\bigg( \dfrac{2}{3}  \biggr)^{-3}$
%--------------------------------------------
 \item $\dfrac{2^{3}}{3^{3}}$
 \item $2^{-4}$
 \item $2^{10}$
 \item $2^{3}\cdot3^{3}$
 \item $2^{6}$
 \item $\bigg( \dfrac{3}{2}  \biggr)^{3}$
\end{multicols}
\end{itemize}








\section*{Άσκηση 2  \hfill \small{30 μονάδες}}
Επιλέξτε τη σωστή απάντηση:
\begin{enumerate}[1)]
 \item το $(6)^{1}$ ισούται με %1
\begin{multicols}{3}
\begin{enumerate}[i)]
 \item $6$
 \item $1$
 \item $0$
\end{enumerate}
\end{multicols}
%--------------------------------------
 \item το  $7^{0}$ ισούται με %t2
\begin{multicols}{3}
 \begin{enumerate}[i)]
 \item $7$
 \item $1$
 \item $0$
 \end{enumerate}
\end{multicols}
 %-----------------------------------
  \item το $8^{-1}$ ισούται με %3
\begin{multicols}{3}
\begin{enumerate}[i)]
 \item $8$
 \item $-1$
 \item $\frac{1}{8}$
\end{enumerate}
\end{multicols}
%------------------------------
 \item το $(9)^{2}\cdot9^{3}$ ισούται με %4
\begin{multicols}{3}
\begin{enumerate}[i)]
 \item $9^{6}$
 \item $9^{5}$
 \item $9^{1}$
\end{enumerate}
\end{multicols}
%------------------------------------
 \item το $(10)^{4}\cdot10^{-12}$ ισούται με %5
\begin{multicols}{3}
\begin{enumerate}[i)]
 \item $10^{16}$
 \item $10^{-8}$
 \item $10^{16}$
\end{enumerate}
\end{multicols}
%---------------------------------------
 \item το $6^{7}:6^{3}$ ισούται με %6
\begin{multicols}{3}
\begin{enumerate}[i)]
 \item $6^{4}$
 \item $6^{10}$
 \item $6^{21}$
\end{enumerate}
\end{multicols}
%-------------------------------------------------
 \item το $5^{4}:5^{-7}$ ισούται με %7
\begin{multicols}{3}
\begin{enumerate}[i)]
 \item $5^{3}$
 \item $5^{11}$
 \item $5^{-3}$
\end{enumerate}
\end{multicols}
%--------------------------------------------
 \item το $9^{-2}: 9^{-5}$ ισούται με %8
\begin{multicols}{3}
\begin{enumerate}[i)]
 \item $9^{3}$
 \item $9^{-7}$
 \item $9^{-3}$
\end{enumerate}
\end{multicols}
%--------------------------------------------
 \item το $(2\cdot6)^{3}$ ισούται με %9
\begin{multicols}{3}
\begin{enumerate}[i)]
 \item $2^{3}\cdot6^{3}$
 \item $2^{9}\cdot6$
 \item $2^{6}\cdot6^{6}$
\end{enumerate}
\end{multicols}
%--------------------------------------------
 \item το $(3\cdot7)^{-4}$ ισούται με %10
\begin{multicols}{3}
\begin{enumerate}[i)]
 \item $3^{-4}\cdot7^{4}$
 \item $3^{4}\cdot7^{-4}$
 \item $3^{-4}\cdot7^{-4}$
\end{enumerate}
\end{multicols}
%--------------------------------------------
 \item το $\biggl( \dfrac{7}{3} \biggr)^{6}$ ισούται με %11
\begin{multicols}{3}
\begin{enumerate}[i)]
 \item $\dfrac{3^{6}}{7^{6}}$
 \item $\dfrac{7^{6}}{3^{6}}$
 \item $\dfrac{7^{-6}}{3^{-6}}$
\end{enumerate}
\end{multicols}
%--------------------------------------------
 \item το $(71^{2})^{3}$ ισούται με %12
\begin{multicols}{3}
\begin{enumerate}[i)]
 \item $71^{1}$
 \item $71^{6}$
 \item $71^{5}$
\end{enumerate}
\end{multicols}
%--------------------------------------------
 \item το $(13^{4})^{3}$ ισούται με %13
\begin{multicols}{3}
\begin{enumerate}[i)]
 \item $13^{12}$
 \item $13^{1}$
 \item $13^{7}$
\end{enumerate}
\end{multicols}
%--------------------------------------------
 \item το $\biggl(\dfrac{5}{6} \biggr)^{-4}$ ισούται με %14
\begin{multicols}{3}
\begin{enumerate}[i)]
 \item $\biggl(\dfrac{5}{6}\biggr)^{4}$
 \item $\biggl(\dfrac{6}{5}\biggr)^{4}$
 \item $\biggl(\dfrac{6}{5}\biggr)^{-4}$
\end{enumerate}
\end{multicols}
%--------------------------------------------
 \item το $\biggl(\dfrac{3}{10}\biggr)^{-2}$ ισούται με %15
\begin{multicols}{3}
\begin{enumerate}[i)]
 \item $\dfrac{10^{-2}}{3^{-2}}$
 \item $\dfrac{3^{2}}{10^{2}}$
 \item $\dfrac{10^{2}}{3^{2}}$
\end{enumerate}
\end{multicols}
%--------------------------------------------
\end{enumerate}


\section*{Άσκηση 3  \hfill \small{} 30 μονάδες}
Αν $x=1$ και $y=2$, να υπολογίσετε τις παρακάτω παραστάσεις:
\begin{enumerate}[i)]
 \item $x^{5}-y^{2}+2xy$
 \item $x^{2}+y^{2}-2xy$
 \item $(x^{2})^{10}+(y^{2})^{3}-5y$
 \item $10\cdot x^{0}+y^{3}+x^{3}\cdot y^{2}$
\end{enumerate}




\section*{Άσκηση 4  \hfill \small{} 30 μονάδες}
Να υπολογίσετε τις παρακάτω παραστάσεις:
\begin{enumerate}[i)]
 \item $3^{0}\cdot 7^{1}+2^{5}-3^{3}$
 \item $(4^{3}:4^{6})\cdot 4^{3}$
 \item $(2^{2})^{7}\cdot 2^{-10}$
 \item $\dfrac{2^{5}\cdot 5^{6}}{(2\cdot 5)^{4}}$
 \item $\dfrac{7^{3}\cdot (11^{2})^{3}}{(-7)^{2}\cdot (7\cdot 11)^{2}\cdot 11^{3}}$
\end{enumerate}



 









%\let\thefootnote\relax\footnotetext{Επιμέλεια : Κασωτάκη Ειρήνη (ikasotaki@gmail.com),    Σμαραγδάκης Κώστας (kesmarag@gmail.com)} 


\end{document}
